\subsection{Weyl Eigenvalue Bounds}

\begin{theorem} \label{thm:weyl-eigenvalue-inequalities-hermitian-sums} \hlt{Weyl Eigenvalue Inequalities}\\
Let $A,B\in\operatorname{Mat}(n,n,\mathbb{C})$ be Hermitian. For $1\leq i\leq n$,
\begin{align*}
\lambda_i(A+B) &\leq \lambda_{i+j}(A)+\lambda_{n-j}(B),
& 0 &\leq j\leq n-i,\\
\lambda_{i-j+1}(A)+\lambda_j(B) &\leq \lambda_i(A+B),
& 1 &\leq j\leq i.
\end{align*}
\begin{proof}
Let $\mathbf{u}_k,\mathbf{v}_k,\mathbf{w}_k$ be increasing eigenbases for $A,A+B,B$.\\
Fix $i$ and $0\leq j\leq n-i$, and set $\mathcal{U}=\operatorname{span}(\mathbf{u}_1,\ldots,\mathbf{u}_{i+j})$, $\mathcal{V}=\operatorname{span}(\mathbf{v}_i,\ldots,\mathbf{v}_n)$, $\mathcal{W}=\operatorname{span}(\mathbf{w}_1,\ldots,\mathbf{w}_{n-j})$.\\
Their dimensions sum to $2n+1$, so choose a unit $\mathbf{x}\in\mathcal{U}\cap\mathcal{V}\cap\mathcal{W}$.\\
Then $\lambda_i(A+B)\leq \mathbf{x}^*(A+B)\mathbf{x} =\mathbf{x}^*A\mathbf{x}+\mathbf{x}^*B\mathbf{x} \leq \lambda_{i+j}(A)+\lambda_{n-j}(B)$.\\
The lower bound follows by applying this upper bound to $-A$ and $-B$ and using $\lambda_k(-C)=-\lambda_{n-k+1}(C)$.\\
\hlt{Structural Commentary}: Weyl bounds come from intersecting the spectral flags of $A$, $B$, and $A+B$.\\
\hlt{Key Insight}: One shared test vector controls all three Rayleigh quotients.
\end{proof}
\end{theorem}

\begin{method} \label{meth:weyl-inequality-method} \hlt{Weyl Bound Method}\\
To bound eigenvalues of $A+B$:
\begin{enumerate}[label=(\roman*)]
\setlength{\itemsep}{0pt}
\item choose the target index $i$;
\item use the upper Weyl family with $j=0,\ldots,n-i$ to get upper bounds;
\item use the lower Weyl family with $j=1,\ldots,i$ to get lower bounds;
\item for positive semidefinite or low-rank $B$, insert the known zero or nonnegative eigenvalues to recover monotonicity or interlacing.
\end{enumerate}
\end{method}
